% PDF-derived LaTeX transcription. Original manuscript source was unavailable.
% Published PDF SHA-256: 674e3165124f54afabda9ca8ed105c30e6630b7fe2af16c529cb89699a842386
% Extracted text SHA-256: 816cfac30641ed1b5e9d29935dffc24a095ddda373cfcb629da41e8e64b6b7d2
\documentclass[11pt]{article}
\usepackage[margin=1in]{geometry}
\usepackage{amsmath,amssymb}
\setlength{\emergencystretch}{3em}
\title{Rotation dependence in an antibody structure tokenizer and sequence-anchored canonicalization}
\author{}
\date{}
\begin{document}
\maketitle
\noindent\textbf{PDF-derived LaTeX transcription.} This editable text follows the PDF extraction. Refer to the published PDF for figures, table layouts, and mathematical notation.

% Extracted chunk 1

Rotation dependence enters a learned antibody
structure tokenizer through its preprocessing and is
removed by sequence-anchored canonicalization
Shodh AI Research

\section*{Abstract}

Learned structure tokenizers turn protein coordinates into per-residue tokens for downstream models, and they are usually assumed to inherit the rotation invariance of their
distance-based encoders. In a preregistered study of a frozen antibody Fv tokenizer, we show
that this assumption can fail through preprocessing alone. The tokenizer’s 160-dimensional
learned features were bit-identical under translation but changed almost completely under rotation (median cosine 0.177 across 976 held-out structures). Exact Shapley attribution traced
99.98\% of the change to Fourier features of coordinates expressed in a per-pose bounding
box. Aligning each antibody to a sequence-anchored joint VH/VL framework frame before
the unchanged encoder made stored tokens bit-identical under every rigid motion while preserving decoding accuracy and sensitivity to CDR-H3 changes. A germline- and geometrybased guard rejected all 976 chain-label swaps and none of 6,976 valid antibodies. On 6,000
sequence-cluster-disjoint antibodies in random poses, the fix restored CDR-H3 structural retrieval (mean average precision 0.165 to 0.531) and VH–VL orientation prediction (Spearman
\ensuremath{\rho} 0.421 to 0.781), although simple invariant descriptors remained stronger for strict loop
retrieval.

Introduction
Protein structure is increasingly passed to learned models as tokens: per-residue vectors or discrete symbols that summarise local and global geometry. Discrete structure alphabets such as
Foldseek’s 3Di1 made structure search as fast as sequence search, and vector-quantised autoencoders2 now supply the structure tokens of multimodal protein language models.3–7 Inversefolding models learn from predicted structures at scale,8 and benchmarks have begun to compare
tokenizers on reconstruction, sensitivity and distinctiveness.9 Antibodies are a prominent application. Fast structure predictors such as ABodyBuilder2 and IgFold10–12 produce Fv models for
millions of sequences from repertoire databases.13, 14 Language models of antibody sequence15, 16
and structure-aware design models17–19 increasingly expect structural inputs in a learned representation.
A structure token should not depend on the arbitrary orientation in which the coordinates happen
to be stored. Geometric deep learning addresses this by building symmetry into the architecture.20
Invariant networks operate on interatomic distances and angles.21, 22 Equivariant networks carry
vector or tensor features that rotate with the input.23–27 Structure-prediction networks reason in
local residue frames.28, 29 Once a tokenizer’s encoder is built from such components, its outputs
are commonly assumed to inherit the encoder’s invariance.
That assumption holds only if nothing else in the pipeline sees raw coordinates, and practical
tokenizers face a tension between invariance and invertibility. A token that can be decoded back
to coordinates must carry some notion of position. Materials representations have confronted
the same trade-off for periodic crystals. Crystal graphs are invariant but not invertible.30, 31
Generative models such as CDVAE represent structures through fractional coordinates and a
1


\end{document}
